% Generated draft: requires mathematical and source review.
\documentclass[11pt]{article}
\usepackage[margin=25mm]{geometry}
\usepackage{fontspec}
\setmainfont{Latin Modern Roman}
\usepackage{amsmath,amssymb,mathtools,unicode-math}
\setmathfont{Latin Modern Math}
\usepackage{enumitem,needspace,xurl,hyperref}
\setlength{\parindent}{0pt}
\setlength{\parskip}{5pt}
\setlength{\emergencystretch}{3em}
\newcommand{\N}{\mathbb{N}}
\newcommand{\Z}{\mathbb{Z}}
\newcommand{\Q}{\mathbb{Q}}
\newcommand{\R}{\mathbb{R}}
\newcommand{\C}{\mathbb{C}}
\newcommand{\F}{\mathbb{F}}
\newcommand{\A}{\mathbb{A}}
\newcommand{\PP}{\mathbb{P}}
\newcommand{\E}{\mathbb{E}}
\newcommand{\eps}{\varepsilon}
\newcommand{\dd}{\,\mathrm{d}}
\begin{document}
\section*{UnsolvedMath ID: 30002774}
\textbf{Source problem code:} catalogue entry
\subsection{Definitions and mathematical statement}
A polynomial \(p(x)\) is real-rooted if every \(z\in\C\) satisfying \(p(z)=0\) belongs to \(\R\).

The supplied record asks for a characterization of the triangulations \(L\) for which the generalized Tchebyshev polynomial \(T_n^L(x)\) is real-rooted for every \(n\). The record does not define the admissible triangulations, the construction of \(T_n^L\), or the range of \(n\). Thus, let \(\mathcal A\) and \(I\) denote, provisionally, the admissible class and full index set intended by that construction. The requested characterization is of
\[
\left\{L\in\mathcal A:\ \forall n\in I\ \forall z\in\C,\quad T_n^L(z)=0\Longrightarrow z\in\R\right\}.
\]
This is a conditional formulation: specifying \(\mathcal A\), \(I\), and the polynomial construction from an authoritative source is necessary for a complete mathematical statement. No particular interpretation of these missing definitions is asserted here.

The supplied literature assessment dated 2026-08-21 reports real-rootedness for important subclasses, including all valid subdivisions of a two-dimensional simplex, and a counterexample arising from a simple three-dimensional subdivision. It reports that a full characterization remains open. These are supplied curation claims, not independently verified results; the meaning of ``valid'' and the counterexample construction are not provided.
\subsection{Short English statement}
Characterize the triangulations \(L\) whose associated generalized Tchebyshev polynomials \(T_n^L(x)\) have only real roots for every index \(n\). The supplied record omits the polynomial construction, admissibility conditions on \(L\), and index range, so the formulation remains conditional on those definitions.
\subsection{Sources}
\begin{itemize}
\item Supplied local UnsolvedMath registry record, permanent ID \(30002774\), source problem code \texttt{catalogue entry}, in the Oberwolfach Reports Open Problems collection. The record gives proposed year 2015 and identifies no individual proposer.
\item \emph{Geometric and Algebraic Combinatorics} (2015), Oberwolfach Reports, pp.\ 285--368, DOI 10.4171/owr/2015/5: \url{https://doi.org/10.4171/owr/2015/5}. This is the supplied workshop-report provenance and also a supplied literature-review source; no exact problem page is given.
\item Hetyei--Nevo, JCTA, supplied literature-source label, DOI 10.1016/j.jcta.2015.08.004: \url{https://doi.org/10.1016/j.jcta.2015.08.004}. No article title or precise result locator is supplied.
\item Literature assessment checked at \texttt{2026-08-21}: supplied dated curation, not independently verified. It labels the problem open, reports real-rootedness for important subclasses including all valid subdivisions of a two-dimensional simplex, reports a simple three-dimensional subdivision counterexample, and states that a full characterization remains open. The record explicitly calls for independent verification of time-sensitive status and cited claims.
\end{itemize}
\end{document}
